\documentclass{article}
\usepackage[utf8]{inputenc}
\usepackage[T1]{fontenc}
\usepackage{graphicx}
\usepackage{hyperref}
\usepackage{geometry}
\geometry{margin=1in}

\title{Britannica Volume 144850375 Row 546}
\date{Encyclopaedia Britannica (1771)}

\begin{document}
\maketitle

\section*{Historical source}
\textbf{Entry ID:} \texttt{entry-britannica-volume-144850375-row-546-546} \\
\textbf{Volume ID:} 144850375 \\
\textbf{Character Count:} 740

\subsection*{Extracted Text}
This means that the number of users are available to the same as they do not have a few times. The number of users can be used in the case where the number of users are able to use the number of users. However, the number of users are also used to generate the number of users and the number of users are able to use the number of users. For example, the number of users are used to compute the number of users and the number of users are able to use the number of users. As mentioned, the number of users are used to compute the number of users and the number of users are able to use the number of users. In addition, the number of users are used to compute the number of users and the number of users are able to use the number of users.

\subsection*{Matched Wikipedia Article}
\textbf{Title:} Arithmetic mean \\
\textbf{Page ID:} 612 \\
\textbf{Revision ID:} 1346126933 \\
\textbf{URL:} \url{https://en.wikipedia.org/wiki/Arithmetic\_mean}


\section*{Source record}
\begin{itemize}
\item \textbf{Dataset:} \texttt{NationalLibraryOfScotland/encyclopaedia\_britannica\_illustrated}
\item \textbf{Dataset Revision:} \texttt{309f765b1ecc83e0a7f7762e44b4a927a7555443}
\item \textbf{Page ID:} \texttt{britannica:309f765b1ecc83e0a7f7762e44b4a927a7555443:train:546} \\
  \textbf{Filename:} \texttt{page-0546-7c6f39d92d2b.jpg} \\
  \textbf{Row Index:} 546 \\
  \textbf{SHA256:} \texttt{7c6f39d92d2b138c...}
\end{itemize}

\begin{figure}[h]
\centering
\includegraphics[width=0.85\textwidth]{../scans/page-0546-7c6f39d92d2b.jpg}
\caption*{Source scan: \texttt{page-0546-7c6f39d92d2b.jpg}}
\end{figure}

\end{document}
